\section{The main window}

The window is organised in five columns, left to right.

\subsection{Observation, site and target}
\label{sec:obs-site-target}

Site selector (built-in observatories merged with your persistent
\code{etc\_sites.json}; \textsf{Save current site} and \textsf{Import site
file} manage it), geographic coordinates and elevation, date and UT time of
the planned observation, the observing-interval selector, the time zone
(IANA name with DST, e.g.\ \code{Europe/Rome}, or a fixed UTC offset), and
the target: RA/Dec fields (decimal degrees or \code{hh:mm:ss} /
\code{dd:mm:ss}) with the optional SIMBAD \textsf{Resolve name} button.
Running the ETC itself never uses the network: only the coordinates in the
fields count.

The OBSERVATORY panel also hosts the \textbf{horizon-profile generator}.
Enter a search radius (1--100, default 10, in km or miles via the unit
selector) and press \textsf{Generate horizon profile}: the program
downloads the Copernicus GLO-30 elevation model around your coordinates
(cached locally, so only the first run of an area needs the internet),
and traces the apparent terrain horizon --- mountains, hills, the
curvature of the Earth and standard refraction included --- in 360
azimuth steps. The computation runs in the background and reports the
highest obstruction when done; the profile is \emph{automatically} saved
both as a CSV and as a PNG plot in \code{data/horizons/} (format in
Sect.~\ref{sec:file-horizon}). The horizon elevation is constrained to the
physical range $-20^\circ$ (the view down from a very high summit) to
$+90^\circ$. \textsf{Display horizon} plots the saved profile for the
current coordinates at any time. When a saved horizon exists for the
current site it is also drawn on the scientific plots: as the terrain
outline on the sky path, and as a shaded exclusion band on the
altitude-versus-time panel, so the target's visibility is shown against
the \emph{local} horizon and not merely the geometric one. The generator needs two
extra libraries (\code{pip install rasterio pyproj}); everything else in
the program works without them.

\subsection{Telescope, detector and conditions}

Telescope diameter, central obstruction, focal length and scalar optics
throughput; optional calibrated instrument-response and slit
resolving-power curves (two-column files with an explicit wavelength unit);
detector pixel size, gain, full well, ADC bit depth, read noise and dark
current; the \textbf{sensor type} (\code{mono}, or \code{osc} with the
R/G/B channel: aperture rates and pixel counts scale by the Bayer fill
fraction, saturation stays per physical pixel; use the channel's effective
QE curve); a \textbf{CMOS gain table} (a text file with one row per gain
setting: setting, e$^-$/ADU, read noise, full well — published by camera
makers and measured on PhotonsToPhotos): pick the setting and the three
detector fields fill consistently; seeing FWHM; the sky-background mode
(\code{ing} = physical night/twilight/day model, \code{fixed\_ab} = your
own observed sky surface brightness, \code{sqm} = your \textbf{zenith SQM
reading} in $V$ mag/arcsec$^2$, with a \textbf{Bortle class} preset that
fills a typical value; band colours and the Moon model are applied on
top); and the \textbf{PSF model} (\code{gaussian} or \code{moffat} with
its $\beta$; Moffat reproduces the wings of real star images and is
recommended); an optional \textbf{seeing-scaling} checkbox that makes the
seeing vary with wavelength and airmass (Kolmogorov turbulence: blue light
more blurred than red, and worse as the target sets --- with it ticked the
value you enter is the zenith $V$ seeing); an optional \textbf{extra
background} in e$^-$/s/pixel, a catch-all for detector glow, stray light or
ghosting that raises the background noise and the peak-pixel prediction;
a \textbf{guiding rms} in arcsec (0 = perfect tracking) that adds
image-motion blur in quadrature to the seeing --- enter your autoguider's
rms error, or an estimate of the unguided drift per exposure; a
\textbf{telluric bands} checkbox that adds the O$_2$/H$_2$O molecular
absorption bands of the red and near-infrared (the deep A band at
7605\,\AA{} in particular) on top of the smooth extinction curve; and the
\textbf{sky annulus pixels} field: the number of pixels you actually use
to measure the background around the star (0 = assume it perfectly
known). A finite value adds the classical $(1+n_{\rm pix}/n_{\rm sky})$
sky-subtraction noise to every background term --- with a small annulus
this is tens of percent of the background noise, so faint-target
predictions become noticeably more honest.
\textsf{Save/Load instrument profile} stores the whole column as a portable
JSON with its QE and response curves.

\subsection{Calculation settings}

Photometry or spectroscopy; exposure time; optional \textbf{target S/N}
(when set, the ETC returns the required exposure instead, flagging
saturation — the solver includes scintillation, so bright-star requests
are honest); the \textbf{stack sub-exposure} preference (0 = automatic)
feeding the stack planner, whose result appears in the label below after
each run. The photometry panel adds the aperture radius, the
\textbf{source geometry}: \code{point}, or \code{extended} with the source
area in arcsec$^2$ for galaxies, nebulae and planets (the magnitude you
enter stays the \emph{total} integrated magnitude), and the optional
\textbf{comparison-star magnitude} for differential photometry — the
mmag-per-frame precision appears in the label below. The spectroscopy
panel holds the resolving power, slit width, extraction height, sampling,
the slitless parameters (Sect.~\ref{sec:spec-guide}), the wavelength range
and the S/N reference wavelength, the \textbf{slit orientation}
(\code{parallactic} or \code{fixed}), and the \textbf{spectrograph
geometry helper}: enter grating lines/mm and collimator/camera focal
lengths and press \textsf{Compute R from geometry} — the Littrow grating
equation fills the R field (which stays editable) and reports whether you
are slit- or seeing-limited. The optional \textbf{clamp} checkbox below
it goes one step further: with it ticked, the engine itself limits the
entered R to what the geometry can deliver, so a typo like R $=$ 100\,000
on a 600 l/mm grating cannot silently inflate every prediction.

\subsection{Filter and template selection}

Two independent filters: the \textbf{reference} filter, in which your
magnitude is expressed (Vega or AB), and the \textbf{observing} filter,
physically in the beam for both photometry and spectroscopy (spectroscopy
defaults to \code{BLANK}, the unfiltered response). Besides the normal
photometric filters (each a \code{_Vega}/\code{_AB} profile pair), the
selector also lists single-file \emph{instrumental} passbands that have no
magnitude of their own and are calibrated on the AB scale: \code{BLANK}
(unfiltered), \code{DARK} (opaque) and neutral-density filters
\code{ND25}, \code{ND50}, … Drop their \code{.xml} into the filters
directory and list them in \code{filters.list} to have them appear. The panel previews the
observing transmission curve with its pivot wavelength, width and zero
point; the template selector searches the catalogue by name and previews
the selected flux distribution over the fixed \SIrange{3000}{9000}{\angstrom}
optical window (so every template is directly comparable); \textsf{Validate
V and B$-$V} checks the template's synthetic colours against the catalogue.
The \textbf{radial velocity} and \textbf{E(B$-$V)} that transform the
template are entered in the \emph{target} box (first column,
Sect.~\ref{sec:obs-site-target}), because they are properties of the
target even though they act on
the template: the radial velocity (km/s, positive receding) shifts all
wavelengths by $(1+v/c)$, and E(B$-$V) reddens the template with the
standard CCM89 interstellar extinction law ($R_V = 3.1$). The magnitude
you enter always stays the \emph{observed} one; the transforms only
change the template's colours, exactly as reddening and motion do to a
real star. After a run the outputs tab reports the resulting synthetic
$(B-V)_0$ of the unreddened template and the reddened $(B-V)$, and — for
spectroscopy — the RV shift in pixels and the velocity sampling
(km/s per pixel) at 4000, 5000, 6000 and \SI{7000}{\angstrom}.

\subsection{Run actions}

\textsf{Run ETC}, the live observatory status (Sun/Moon altitude, current
airmass of the target), and \textsf{Exit} (which saves the configuration).
